%----------------------------------------------------------------------------------------
% NOTATIE EN FIGUREN VAN DE CURSUS REËLE GETALLEN (2de graad)
% Staat in de preambule en niet in een hoofdstuk: ook de lijsten met oplossingen gebruiken dit.
% Vertrokken van notatie.tex van Integralen; wat hier niet nodig is, is geschrapt.
%----------------------------------------------------------------------------------------
% Bewaarde grafiek, geplaatst door maak-pdf.py: bestand en pagina in figuren-cache/
\newcommand{\grafiekbestand}[2]{\includegraphics[page=#2,trim=3bp 3bp 3bp 3bp]{figuren-cache/#1}}

% Getallenverzamelingen
\providecommand{\N}{\mathbb{N}}
\providecommand{\Z}{\mathbb{Z}}
\providecommand{\Q}{\mathbb{Q}}
\providecommand{\R}{\mathbb{R}}
% Repeterend decimaal getal: \rep{0,}{3} geeft 0,333... met een streep boven het blok dat herhaalt
\newcommand{\rep}[2]{#1\overline{#2}}

% Zoals stappen, maar met een bredere kolom voor de berekening: voor lange formules
\newenvironment{stappenbreed}{%
  \par\renewcommand{\arraystretch}{1.25}%
  \noindent\begin{tabular}{@{}>{$\displaystyle}p{9.8cm}<{$}@{\hspace{4mm}}>{\RaggedRight\small\itshape\arraybackslash}p{4.9cm}@{}}%
}{\end{tabular}\par}

% Zoals stappen, met eigen kolombreedtes: \begin{stappenvrij}{6.2cm}{3.8cm}. Voor naast een figuur.
\newenvironment{stappenvrij}[2]{%
  \par\renewcommand{\arraystretch}{1.25}%
  \noindent\begin{tabular}{@{}>{$\displaystyle}p{#1}<{$}@{\hspace{3mm}}>{\RaggedRight\small\itshape\arraybackslash}p{#2}@{}}%
}{\end{tabular}\par}

% Tekst of berekening links, figuur rechts: \naast{links}{rechts}
\newsavebox{\naastbox}
\newcommand{\naast}[3][]{%
  \par\noindent\sbox{\naastbox}{#3}%
  \begin{minipage}[t]{\dimexpr\linewidth-\wd\naastbox-3mm\relax}\vspace{0pt}#2\end{minipage}\hfill
  \begin{minipage}[t]{\wd\naastbox}\vspace{0pt}\usebox{\naastbox}\end{minipage}\par\addvspace{1mm}}
% Een rij figuren, gecentreerd
\newcommand{\figuren}[1]{\par\addvspace{1mm}{\centering #1\par}\addvspace{1.5mm}}
% Opgave met een figuur rechts: \begin{metgrafiek}{tikz} opgave en deelvragen \end{metgrafiek}
\newsavebox{\grafiekbox}
\newenvironment{metgrafiek}[2][]{%
  \sbox{\grafiekbox}{#2}%
  \begin{minipage}[t]{\dimexpr\linewidth-\wd\grafiekbox-3mm\relax}%
}{%
  \end{minipage}\hfill
  \begin{minipage}[t]{\wd\grafiekbox}\vspace{-2mm}\usebox{\grafiekbox}\end{minipage}\par}

% GETALLENAS, binnen een tikzpicture. Zet de schaal met x=... (bv. [x=1cm] of [x=.5cm]).
% \getallenas{min}{max}: as met een pijl, streepjes en getallen bij de gehele getallen van min tot max.
% \getallenasleeg{min}{max}: hetzelfde zonder getallen (voor eigen labels).
% \getallenasstap{min}{stap}{max}: streepjes en getallen met een andere stap, bv. {0}{0.1}{1} (de getallen met komma).
\newcommand{\getallenasleeg}[2]{%
  \draw[->,grijs] ({#1-.45},0) -- ({#2+.55},0);
  \foreach \t in {#1,...,#2}{\draw[grijs] (\t,-.09) -- (\t,.09);}}
\newcommand{\getallenas}[2]{%
  \getallenasleeg{#1}{#2}
  \foreach \t in {#1,...,#2}{\node[below,inner sep=2pt,text=black,font=\scriptsize] at (\t,-.08) {$\t$};}}
\newcommand{\getallenasstap}[3]{%
  \draw[->,grijs] ({#1-.45*(#2)*10},0) -- ({#3+.55*(#2)*10},0);
  \pgfmathsetmacro{\gtwee}{#1+#2}%
  \foreach \t in {#1,\gtwee,...,#3}{\draw[grijs] (\t,-.09) -- (\t,.09);}}
% Een punt op de as, met een label erboven: \aspunt{x}{label} of \aspunt[below]{x}{label}
\newcommand{\aspunt}[3][above]{\fill[accent] (#2,0) circle (2.2pt) node[#1,inner sep=3pt,text=black,font=\footnotesize] {$#3$};}
% Absolute waarde: \abs{-3} geeft |-3| met een minteken (niet een aftrekking) na de streep; \Abs met grotere strepen
\newcommand{\abs}[1]{\lvert #1\rvert}
\newcommand{\Abs}[1]{\bigl\lvert #1\bigr\rvert}
